% =========================================================
% Algebra of completeness bounds
% =========================================================
\subsection{Algebraic Properties of Supremum and Infimum}

% =========================================================
% Uniqueness of the Maximum
% =========================================================
\begin{proposition}[Uniqueness of the Maximum]
\label{prop:uniqueness-of-the-maximum}
If a set $A \subseteq \mathbb{R}$ has a maximum, then that maximum is unique.

\smallskip
\noindent
\hyperref[prf:uniqueness-of-the-maximum]{\textit{Go to proof.}}
\end{proposition}

\begin{remark*}[Standard quantified statement]
\[
\bigl(\operatorname{Maximum}(m_1, A) \land \operatorname{Maximum}(m_2, A)\bigr) \implies m_1 = m_2
\]
\end{remark*}

\begin{remark*}[Interpretation]
The maximum is the "greatest" element. Since the order $\le$ on $\mathbb{R}$ is antisymmetric, a set cannot have two distinct elements that are both greater than or equal to all other elements.
\end{remark*}

% =========================================================
% Uniqueness of the Supremum
% =========================================================
\begin{proposition}[Uniqueness of the Supremum]
\label{prop:uniqueness-of-the-supremum}
If a set $A \subseteq \mathbb{R}$ has a supremum, then that supremum is unique.

\smallskip
\noindent
\hyperref[prf:uniqueness-of-the-supremum]{\textit{Go to proof.}}
\end{proposition}

\begin{remark*}[Standard quantified statement]
\[
\bigl(\operatorname{Supremum}(s_1, A) \land \operatorname{Supremum}(s_2, A)\bigr) \implies s_1 = s_2
\]
\end{remark*}

\begin{remark*}[Interpretation]
While there are many guards (upper bounds), there is only one "least" guard. This follows from the property that the supremum is the minimum of the set of upper bounds.
\end{remark*}

% =========================================================
% Maximum Implies Supremum
% =========================================================
\begin{proposition}[Maximum Implies Supremum]
\label{prop:maximum-implies-supremum}
If a set $A \subseteq \mathbb{R}$ has a maximum $m$, then $m$ is also the supremum of $A$.

\smallskip
\noindent
\hyperref[prf:maximum-implies-supremum]{\textit{Go to proof.}}
\end{proposition}

\begin{remark*}[Standard quantified statement]
\[
\operatorname{Maximum}(m, A) \implies \operatorname{Supremum}(m, A)
\]
\end{remark*}

% =========================================================
% Supremum in the Set Is the Maximum
% =========================================================
\begin{proposition}[Supremum in the Set Is the Maximum]
\label{prop:supremum-in-the-set-is-the-maximum}
Let $s = \sup A$. If $s \in A$, then $s$ is the maximum of $A$.

\smallskip
\noindent
\hyperref[prf:supremum-in-the-set-is-the-maximum]{\textit{Go to proof.}}
\end{proposition}

\begin{remark*}[Standard quantified statement]
\[
\bigl(\operatorname{Supremum}(s, A) \land (s \in A)\bigr) \implies \operatorname{Maximum}(s, A)
\]
\end{remark*}



\input{volume-iii/analysis/bounding/notes/bounds-algebra/figure-supremum-algebra}



% =========================================================
% Subset Inclusion Preserves Upper Bounds
% =========================================================
\begin{lemma}[Subset Inclusion Preserves Upper Bounds]
\label{lem:subset-inclusion-preserves-upper-bounds}
Let $A \subseteq B \subseteq \mathbb{R}$. If $u$ is an upper bound of $B$, then $u$ is an upper bound of $A$.

\smallskip
\noindent
\hyperref[prf:subset-inclusion-preserves-upper-bounds]{\textit{Go to proof.}}
\end{lemma}

\begin{remark*}[Standard quantified statement]
\[
\bigl(\operatorname{Subset}(A, B) \land \operatorname{UpperBound}(u, B)\bigr) \implies \operatorname{UpperBound}(u, A)
\]
\end{remark*}

% =========================================================
% Supremum Is Monotone Under Inclusion
% =========================================================
\begin{theorem}[Supremum Is Monotone Under Inclusion]
\label{thm:supremum-is-monotone-under-inclusion}
Let $A$ and $B$ be non-empty subsets of $\mathbb{R}$ that are bounded above. If $A \subseteq B$, then $\sup A \le \sup B$.

\smallskip
\noindent
\hyperref[prf:supremum-is-monotone-under-inclusion]{\textit{Go to proof.}}
\end{theorem}

\begin{remark*}[Standard quantified statement]
\[
\bigl(\operatorname{Subset}(A, B) \land \operatorname{Supremum}(s_A, A) \land \operatorname{Supremum}(s_B, B)\bigr) \implies s_A \le s_B
\]
\end{remark*}

% =========================================================
% Upper Bound Comparison with the Supremum
% =========================================================
\begin{proposition}[Upper Bound Comparison with the Supremum]
\label{prop:upper-bound-comparison-with-the-supremum}
Let $s = \sup A$. For any $u \in \mathbb{R}$, $u$ is an upper bound of $A$ if and only if $s \le u$.

\smallskip
\noindent
\hyperref[prf:upper-bound-comparison-with-the-supremum]{\textit{Go to proof.}}
\end{proposition}

\begin{remark*}[Standard quantified statement]
\[
\operatorname{Supremum}(s, A) \implies \forall u \in \mathbb{R} \; \bigl(\operatorname{UpperBound}(u, A) \iff s \le u \bigr)
\]
\end{remark*}

% =========================================================
% Every Element Lies Below the Supremum
% =========================================================
\begin{proposition}[Every Element Lies Below the Supremum]
\label{prop:every-element-lies-below-the-supremum}
If $s = \sup A$, then for every $x \in A$, $x \le s$.

\smallskip
\noindent
\hyperref[prf:every-element-lies-below-the-supremum]{\textit{Go to proof.}}
\end{proposition}

\begin{remark*}[Standard quantified statement]
\[
\operatorname{Supremum}(s, A) \implies \forall x \in A \; (x \le s)
\]
\end{remark*}

% =========================================================
% Infimum Is Less Than or Equal to the Supremum
% =========================================================
\begin{theorem}[Infimum Is Less Than or Equal to the Supremum]
\label{thm:infimum-less-than-supremum}
For any non-empty bounded set $A \subseteq \mathbb{R}$, $\inf A \le \sup A$.

\smallskip
\noindent
\hyperref[prf:infimum-less-than-supremum]{\textit{Go to proof.}}
\end{theorem}

\begin{remark*}[Standard quantified statement]
\[
\bigl(\operatorname{Infimum}(t, A) \land \operatorname{Supremum}(s, A)\bigr) \implies t \le s
\]
\end{remark*}

% =========================================================
% Order Comparison of Suprema
% =========================================================
\begin{proposition}[Order Comparison of Suprema]
\label{prop:order-comparison-of-suprema}
Let $A, B \subseteq \mathbb{R}$ be non-empty and bounded above. If for every $x \in A$ there exists $y \in B$ such that $x \le y$, then $\sup A \le \sup B$.

\smallskip
\noindent
\hyperref[prf:order-comparison-of-supremum]{\textit{Go to proof.}}
\end{proposition}

\begin{remark*}[Standard quantified statement]
\[
\bigl(\forall x \in A \; \exists y \in B \; (x \le y)\bigr) \implies \sup A \le \sup B
\]
\end{remark*}




% --- Algebra of Supremum Cluster ---

\begin{theorem}[Translation Invariance of the Supremum]
\label{thm:translation-invariance-supremum}
Let $A \subseteq \mathbb{R}$ be a non-empty set bounded above, and let $c \in \mathbb{R}$. Define the translated set $A + c = \{a + c : a \in A\}$. Then $\sup(A + c) = \sup A + c$.

\smallskip
\noindent
\hyperref[prf:translation-invariance-supremum]{\textit{Go to proof.}}
\end{theorem}

\begin{theorem}[Scalar Multiplication by a Positive Scalar]
\label{thm:scalar-mult-positive}
Let $A \subseteq \mathbb{R}$ be a non-empty set bounded above, and let $k \in \mathbb{R}$ such that $k > 0$. Define the scaled set $kA = \{ka : a \in A\}$. Then $\sup(kA) = k \sup A$.

\smallskip
\noindent
\hyperref[prf:scalar-mult-positive]{\textit{Go to proof.}}
\end{theorem}

\begin{theorem}[Scalar Multiplication by a Negative Scalar]
\label{thm:scalar-mult-negative}
Let $A \subseteq \mathbb{R}$ be a non-empty set bounded below, and let $k \in \mathbb{R}$ such that $k < 0$. Define the scaled set $kA = \{ka : a \in A\}$. Then $\sup(kA) = k \inf A$.

\smallskip
\noindent
\hyperref[prf:scalar-mult-negative]{\textit{Go to proof.}}
\end{theorem}

\begin{theorem}[Negation Exchanges Supremum and Infimum]
\label{thm:negation-exchange-sup-inf}
Let $A \subseteq \mathbb{R}$ be a non-empty set. If $A$ is bounded below, then the set $-A = \{-a : a \in A\}$ is bounded above and $\sup(-A) = -\inf A$.

\smallskip
\noindent
\hyperref[prf:negation-exchange-sup-inf]{\textit{Go to proof.}}
\end{theorem}

\begin{remark*}[Standard quantified statement]
\[
\forall A \subseteq \mathbb{R} \; \bigl( \operatorname{NonemptySet}(A) \land \operatorname{BoundedBelow}(A) \bigr) \implies \sup(-A) = -\inf A
\]
\end{remark*}

\begin{remark*}[Definition predicate reading]
\[
\operatorname{Supremum}(-\inf A, -A)
\]
\end{remark*}

\begin{remark*}[Interpretation]
This theorem establishes that negation is an order-reversing map that swaps extremal boundaries. Geometrically, reflecting a set across the origin transforms its "floor" into a "ceiling." This allows us to prove statements about infima by transforming them into statements about suprema, which are often more intuitive to manipulate in the real number system.
\end{remark*}




\begin{theorem}[Supremum of a Sum Set]
\label{thm:supremum-sum-set}
Let $A, B \subseteq \mathbb{R}$ be non-empty sets bounded above. Define the Minkowski sum $A + B = \{a + b : a \in A, b \in B\}$. Then $\sup(A + B) = \sup A + \sup B$.

\smallskip
\noindent
\hyperref[prf:supremum-sum-set]{\textit{Go to proof.}}
\end{theorem}

\begin{remark*}[Standard quantified statement]
\[
\begin{aligned}
&\forall A, B \subseteq \mathbb{R} : \\
&\quad \bigl( \operatorname{NonemptySet}(A) \land \operatorname{NonemptySet}(B) \land \operatorname{BoundedAbove}(A) \land \operatorname{BoundedAbove}(B) \bigr) \\
&\quad \implies \sup(A + B) = \sup A + \sup B
\end{aligned}
\]
\end{remark*}

\begin{remark*}[Interpretation]
This theorem confirms that the supremum operator behaves linearly with respect to set addition. It implies that the "ceiling" of a combined system is simply the sum of the individual ceilings. This allows bounding the sum of two independent stochastic displacements, the total bound, by simply summing the individual Lipschitz or supremum constraints.
\end{remark*}





















\begin{theorem}[Supremum of a Difference Set]
\label{thm:supremum-difference-set}
Let $A, B \subseteq \mathbb{R}$ be non-empty sets such that $A$ is bounded above and $B$ is bounded below. Define $A - B = \{a - b : a \in A, b \in B\}$. Then $\sup(A - B) = \sup A - \inf B$.

\smallskip
\noindent
\hyperref[prf:supremum-difference-set]{\textit{Go to proof.}}
\end{theorem}

\begin{theorem}[Supremum of a Dilation]
\label{thm:supremum-dilation}
Let $A \subseteq \mathbb{R}$ be a non-empty bounded set and let $\lambda \in \mathbb{R}$. The dilation $\lambda A$ satisfies $\sup(\lambda A) = \max\{\lambda \sup A, \lambda \inf A\}$.

\smallskip
\noindent
\hyperref[prf:supremum-dilation]{\textit{Go to proof.}}
\end{theorem}

\begin{theorem}[Supremum of the Absolute-Value Image]
\label{thm:supremum-absolute-value-image}
Let $A \subseteq \mathbb{R}$ be a non-empty bounded set. Define $|A| = \{|a| : a \in A\}$. Then $\sup |A| = \max\{|\sup A|, |\inf A|\}$.

\smallskip
\noindent
\hyperref[prf:supremum-absolute-value-image]{\textit{Go to proof.}}
\end{theorem}

\begin{proposition}[Bounds for Suprema and Infima Under Set Addition]
\label{prop:bounds-sum-set}
Let $A, B \subseteq \mathbb{R}$ be non-empty bounded sets. For the sum set $A+B$, the following inequalities hold: $\inf A + \inf B \le \inf(A + B)$ and $\sup(A + B) \le \sup A + \sup B$.

\smallskip
\noindent
\hyperref[prf:bounds-sum-set]{\textit{Go to proof.}}
\end{proposition}



% --- Containing-set / Ambient-set Cluster ---

\begin{proposition}[Upper Bounds Depend on the Ambient Order]
\label{prop:upper-bounds-ambient-order}
Let $A \subseteq S$. The property of $u \in S$ being an upper bound of $A$ is defined relative to the ordered set $(S, \le)$. Specifically, $\operatorname{UpperBound}(u, A)$ holds if $\forall x \in A, x \le u$ in $S$.

\smallskip
\noindent
\hyperref[prf:upper-bounds-ambient-order]{\textit{Go to proof.}}
\end{proposition}

\begin{proposition}[Suprema Depend on the Ambient Set]
\label{prop:suprema-ambient-set}
Let $A \subseteq S \subseteq S'$. A value $s$ may satisfy $\operatorname{Supremum}(s, A)$ relative to $S$ but fail to satisfy it relative to $S'$, or vice versa, depending on the existence of elements in the ambient set.

\smallskip
\noindent
\hyperref[prf:suprema-ambient-set]{\textit{Go to proof.}}
\end{proposition}

\begin{theorem}[Ambient Existence of Supremum]
\label{thm:ambient-existence-supremum}
A set $A$ may have a supremum in an ordered set $S'$ but lack a supremum in a subset $S \subset S'$, even if $A \subseteq S$.

\smallskip
\noindent
\hyperref[prf:ambient-existence-supremum]{\textit{Go to proof.}}
\end{theorem}

\begin{lemma}[Rational Gap Example for Suprema]
\label{lem:rational-gap-suprema}
Let $A = \{q \in \mathbb{Q} : q^2 < 2\}$. The set $A$ is non-empty and bounded above in $\mathbb{Q}$, but there exists no $s \in \mathbb{Q}$ such that $\operatorname{Supremum}(s, A)$ relative to the ambient set $\mathbb{Q}$.

\smallskip
\noindent
\hyperref[prf:rational-gap-suprema]{\textit{Go to proof.}}
\end{lemma}




\input{volume-iii/analysis/bounding/notes/bounds-algebra/figure-bounds-algebra}



















